% Momento en que se materializa la identidad de la instancia productora.
% Dos trayectos idénticos que sólo difieren en el estado de la petición.
\begin{singlespace}
\begin{tikzpicture}[figbase]

    % ---------------------------------------------------------------- carril A
    \draw[draw=figGrisSuave, line width=0.5pt] (-1.55,3.35) -- (14.05,3.35);
    \node[figcarril] at (-1.55,3.07) {Petición inicial, sin destino específico};

    \node[fignf]  (ca) at (0,0)    {NF consumidora};
    \node[figpep] (pa) at (3.1,0)  {PEP previo\\al SCP};
    \node[figinfra, minimum height=1.7cm] (sa) at (6.2,0) {};
    \node[text=figInfraTexto, anchor=north, inner sep=5pt] at (sa.north) {SCP};
    \node[draw=figPEPborde, fill=white, text=figPEPtexto,
          line width=0.7pt, dash pattern=on 2.5pt off 1.5pt, rounded corners=1.5pt,
          font=\scriptsize, minimum width=1.75cm, minimum height=0.44cm]
          (sela) at (6.2,-0.42) {selección};
    \node[figpep] (qa) at (9.3,0)  {PEP posterior\\al SCP};
    \node[fignf]  (ua) at (12.4,0) {UDM\_a};

    \node[fignrf] (nrf) at (6.2,2.32) {NRF};

    \draw[figflujo] (ca) -- (pa);
    \draw[figflujo] (pa) -- (sa);
    \draw[figflujo] (sa) -- (qa);
    \draw[figflujo] (qa) -- (ua);
    \draw[figaux, <->] (sa.north) -- (nrf.south)
        node[figrot, midway, right=3pt] {descubrimiento delegado};

    % Línea de vida del discriminador en el carril A.
    \draw[draw=figNegBorde, line width=1pt, dash pattern=on 3.5pt off 2pt]
        (-1.2,-1.52) -- (6.2,-1.52);
    \draw[draw=figPosBorde, line width=1pt] (6.2,-1.52) -- (13.6,-1.52);
    \draw[draw=figInfraBorde, line width=0.8pt] (6.2,-1.36) -- (6.2,-1.68);
    \node[figrot, anchor=south, text=figInfraTexto] at (6.2,-1.34) {se materializa};
    \node[figrot, anchor=north, text=figNegTexto] at (2.5,-1.70)
        {identidad de la instancia productora: aún no materializada};
    \node[figrot, anchor=north, text=figPosTexto] at (9.9,-1.70) {materializada};

    % ---------------------------------------------------------------- carril B
    \draw[draw=figGrisSuave, line width=0.5pt] (-1.55,-2.42) -- (14.05,-2.42);
    \node[figcarril] at (-1.55,-2.70) {Petición posterior, asociada a un recurso};

    \node[fignf]  (cb) at (0,-4.05)    {NF consumidora};
    \node[figpep] (pb) at (3.1,-4.05)  {PEP previo\\al SCP};
    \node[figinfra] (sb) at (6.2,-4.05) {SCP};
    \node[figpep] (qb) at (9.3,-4.05)  {PEP posterior\\al SCP};
    \node[fignf]  (ub) at (12.4,-4.05) {UDM\_a};

    \draw[figflujo, draw=figPosBorde, line width=1.1pt] (cb) -- (pb);
    \draw[figflujo, draw=figPosBorde, line width=1.1pt] (pb) -- (sb);
    \draw[figflujo, draw=figPosBorde, line width=1.1pt] (sb) -- (qb);
    \draw[figflujo, draw=figPosBorde, line width=1.1pt] (qb) -- (ub);

    % La cabecera que transporta el destino se rotula fuera de la banda de nodos.
    \node[figrot, text=figPosTexto] (tag) at (1.5,-3.16)
        {transporta \texttt{3gpp-Sbi-Target-apiRoot}};
    \draw[draw=figGrisSuave, line width=0.5pt, dash pattern=on 1.5pt off 1.5pt]
        (tag.south) -- (1.6,-4.05);

    % Línea de vida del discriminador en el carril B.
    \draw[draw=figPosBorde, line width=1pt] (-1.2,-5.35) -- (13.6,-5.35);
    \node[figrot, anchor=north, text=figPosTexto] at (6.2,-5.53)
        {identidad de la instancia productora: materializada desde que el consumidor emite la petición};

\end{tikzpicture}
\end{singlespace}
